\documentclass[letterpaper,11pt]{article}

\usepackage[utf8]{inputenc}
\usepackage[T1]{fontenc}
\usepackage{graphicx}
\usepackage{latexsym}
\usepackage[empty]{fullpage}
\usepackage{titlesec}
\usepackage{marvosym}
\usepackage[usenames,dvipsnames]{color}
\usepackage{verbatim}
\usepackage{enumitem}
\usepackage[hidelinks]{hyperref}
\usepackage{fancyhdr}
\usepackage{tabularx}
\input{glyphtounicode}

\pagestyle{fancy}
\fancyhf{} 
\fancyfoot{}
\renewcommand{\headrulewidth}{0pt}
\renewcommand{\footrulewidth}{0pt}
\setlength{\footskip}{5pt} 

% Kenar bosluklari
\addtolength{\oddsidemargin}{-0.5in}
\addtolength{\evensidemargin}{-0.5in}
\addtolength{\textwidth}{1in}
\addtolength{\topmargin}{-0.5in}
\addtolength{\textheight}{1.0in}

\urlstyle{same}

\raggedbottom
\raggedright
\setlength{\tabcolsep}{0in}

% Baslik formati
\titleformat{\section}{
  \vspace{-4pt}\scshape\raggedright\large
}{}{0em}{}[\color{black}\titlerule \vspace{-5pt}]

% PDF ATS uyumlulugu icin
\pdfgentounicode=1

% Ozel listeler icin komutlar
\newcommand{\resumeItem}[1]{
  \item\small{
    {#1 \vspace{-2pt}}
  }
}

\newcommand{\resumeSubheading}[4]{
  \vspace{-2pt}\item
    \begin{tabular*}{0.97\textwidth}[t]{l@{\extracolsep{\fill}}r}
      \textbf{#1} & #2 \\
      \textit{\small#3} & \textit{\small #4} \\
    \end{tabular*}\vspace{-7pt}
}

\newcommand{\resumeProjectHeading}[2]{
    \item
    \begin{tabular*}{0.97\textwidth}{l@{\extracolsep{\fill}}r}
      \small#1 & #2 \\
    \end{tabular*}\vspace{-7pt}
}

\newcommand{\resumeSubItem}[1]{\resumeItem{#1}\vspace{-4pt}}
\renewcommand\labelitemii{$\vcenter{\hbox{\tiny$\bullet$}}$}
\newcommand{\resumeSubHeadingListStart}{\begin{itemize}[leftmargin=0.15in, label={}]}
\newcommand{\resumeSubHeadingListEnd}{\end{itemize}}
\newcommand{\resumeItemListStart}{\begin{itemize}}
\newcommand{\resumeItemListEnd}{\end{itemize}\vspace{-5pt}}

\begin{document}

%----------HEADER----------
\begin{minipage}[c]{0.75\textwidth}
    \textbf{\Huge \scshape Sarp Onaran} \\ \vspace{6pt}
    \small 
    0531 234 25 65 $|$ \href{mailto:sarponaran@icloud.com}{\underline{sarponaran@icloud.com}} \\
    \href{https://linkedin.com/in/sarp-onaran}{\underline{linkedin.com/in/sarp-onaran}} $|$ \href{https://github.com/sarp-onaran}{\underline{github.com/sarp-onaran}} \\
    İstanbul, Türkiye
\end{minipage}
\hfill
\begin{minipage}[c]{0.20\textwidth}
    \raggedleft
    \includegraphics[width=2.8cm]{foto.jpeg}
\end{minipage}
\vspace{10pt}

%-----------OBJECTIVE-----------
\section{Hakkımda \& Kariyer Hedefi}
Işık Üniversitesi Yazılım Mühendisliği bölümünde 3. sınıf öğrencisiyim. Nesne yönelimli programlama, veri yapıları ve backend mimarileri üzerine yoğunlaşarak karmaşık finansal ve operasyonel senaryoları çözen projeler geliştirdim. Akademik ve kişisel projelerimde edindiğim teknik derinliği profesyonel süreçlere entegre ederek stajyer rolünde ekibinize değer katmayı hedefliyorum.

%-----------EDUCATION-----------
\section{Eğitim}
  \resumeSubHeadingListStart
    \resumeSubheading
      {Işık Üniversitesi}{İstanbul, Türkiye}
      {Yazılım Mühendisliği (İngilizce) - Lisans}{2019 -- Devam Ediyor}
  \resumeSubHeadingListEnd

%-----------PROJECTS-----------
\section{Projeler \& Akademik Çalışmalar}
    \resumeSubHeadingListStart

      \resumeProjectHeading
          {\textbf{MarketSim PRO} $|$ \emph{JavaScript, HTML/CSS, Chart.js, Vercel} $|$ \href{https://pazar-simulasyonu.vercel.app}{\underline{Canlı Demo}}}{Kişisel Proje}
          \resumeItemListStart
            \resumeItem{Ürünlerin pazar başarısını tahmin eden, 10.000 iterasyonluk Monte Carlo simülasyonu çalıştıran interaktif bir web uygulaması geliştirildi.}
            \resumeItem{World Bank, IMF ve TÜİK verilerine dayalı ülke bazlı ekonomik parametreler entegre edilerek gerçekçi pazar analizi altyapısı kuruldu.}
            \resumeItem{Chart.js ile dinamik veri görselleştirme, tema sistemi (Dark/Light/Ocean) ve Hızlı/Detaylı mod seçenekleri ile kullanıcı deneyimi optimize edildi.}
          \resumeItemListEnd

      \resumeProjectHeading
          {\textbf{FlexPay Smart Wallet Engine} $|$ \emph{Java, SQLite, Maven, Design Patterns}}{Kişisel Proje}
          \resumeItemListStart
            \resumeItem{Kurumsal yemek, kredi kartı ve sadakat puanlarını yöneten; "Küsurat Kumbarası (Micro-Saving)" algoritmasıyla otomatik tasarruf sağlayan akıllı parçalı ödeme (Split-Payment) motoru geliştirildi.}
            \resumeItem{Farklı sektör kurallarını dinamik yönetmek için Strategy, veritabanı bağlantısı için Singleton tasarım kalıpları ve SQLite (JDBC) entegrasyonu uygulandı.}
            \resumeItem{Katmanlı mimari (model/repository/service) üzerine kurulan sistemde, özel exception sınıfları yazılarak tüm işlemlerin zaman damgalı loglanması (Audit Trail) sağlandı.}
          \resumeItemListEnd

      \resumeProjectHeading
          {\textbf{Algorithmic Trading \& Portfolio Management API} $|$ \emph{Java, SQLite, REST API}}{Kişisel Proje}
          \resumeItemListStart
            \resumeItem{CoinGecko API üzerinden anlık piyasa verilerini çekerek çalışan algoritmik bir ticaret botu arka uç sistemi tasarlandı.}
            \resumeItem{SQLite portföy yönetimi, kurumsal hata yönetimi (Exception Handling) ve dosya sistemine loglama (Audit Trail) entegre edilerek profesyonel JavaDoc dökümantasyonu yazıldı.}
          \resumeItemListEnd

      \resumeProjectHeading
          {\textbf{File System Simulation} $|$ \emph{Java, Data Structures}}{Akademik Proje}
          \resumeItemListStart
            \resumeItem{Tree, LinkedList, Stack ve HashTable veri yapıları dış kütüphane kullanılmadan sıfırdan implemente edilerek hiyerarşik dosya sistemi modeli tasarlandı.}
            \resumeItem{Sistem mimarisi 2 kişilik ekip çalışması ile Agile prensiplerine uygun şekilde geliştirildi.}
          \resumeItemListEnd

    \resumeSubHeadingListEnd

%-----------TECHNICAL SKILLS-----------
\section{Yetenekler}
 \begin{itemize}[leftmargin=0.15in, label={}]
    \small{\item{
     \textbf{Programlama Dilleri \& Araçlar}{: Java (Orta Seviye), JavaScript, HTML/CSS, SQL (Başlangıç Seviyesi), Git/GitHub, Temel AWS} \\
     \textbf{Mimari \& Konseptler}{: OOP, Data Structures, Design Patterns (Strategy, Singleton), REST API, Monte Carlo Simülasyonu}
    }}
 \end{itemize}

%-----------LANGUAGES-----------
\section{Yabancı Dil}
 \begin{itemize}[leftmargin=0.15in, label={}]
    \small{\item{
     \textbf{İngilizce}{: B1-B2 (Akademik Okuma/Yazma ve Teknik Konuşma Yetkinliği)}
    }}
 \end{itemize}

%-----------REFERENCES-----------
\section{Referanslar}
 \begin{itemize}[leftmargin=0.15in, label={}]
    \small{\item{
     \textbf{Birtan Atasoy} -- Sistem ve Altyapı Uzmanı, Tanı Pazarlama ve İletişim Hizmetleri A.Ş. (0543 506 97 83)
    }}
 \end{itemize}

\end{document}
